\begin{table}[tbp]
\centering
\caption{Existing robust and clustered benchmark evidence}
\label{tab:bench-vce}
\small
\begin{threeparttable}
\begin{tabular}{llrrrr}
\toprule
$N$ & VCE & Threads & \texttt{fereg} & \texttt{reghdfe} & Ratio \\
\midrule
1m & HC1 robust & 1 & 10.982 & 46.182 & 0.238 \\
1m & HC1 robust & 16 & 3.000 & 14.086 & 0.213 \\
1m & Cluster (one way) & 1 & 6.875 & 25.006 & 0.275 \\
1m & Cluster (one way) & 16 & 4.437 & 30.308 & 0.146 \\
10m & HC1 robust & 1 & 161.455 & 442.199 & 0.365 \\
10m & HC1 robust & 16 & 87.216 & 369.233 & 0.236 \\
10m & Cluster (one way) & 1 & 131.854 & 350.998 & 0.376 \\
10m & Cluster (one way) & 16 & 60.488 & 205.534 & 0.294 \\
\midrule
\multicolumn{5}{l}{Geometric mean across eight cells} & 0.258 \\
\bottomrule
\end{tabular}
\begin{tablenotes}[flushleft]\footnotesize
\item Notes: Seconds are within-shard warm medians for worker--firm models with 30 finite controls. Ratio is \texttt{fereg}/\texttt{reghdfe}. Every shard ran on a 32-processor SCC node. Stata/IC used one processor in the one-thread cells and four processors in the 16-thread cells; the native \texttt{fereg} request was validated separately. See Section~\ref{sec:benchmarks} for scope and qualification cautions.
\end{tablenotes}
\end{threeparttable}
\end{table}
